\section{Allocation on a separated tree: rank data}
\label{sec:allocation}

We first solve an allocation problem on a fixed finite tree.  The allocation
uses conditional distributions of hidden subtree counts.  Convex functions
of their rank probabilities give compatible lower bounds on the amounts
allocated to subtrees.  A small, size-dependent change in the functions
provides the slack needed to move those amounts slowly.

Let $\mathcal T$ be a finite rooted tree whose leaves have depth $J$.
Write $|v|$ for depth and $i\prec v$ when $i$ is a child of $v$.
Fix $r_j=r_0\tau^{-j}$, where $r_0>0$ and $\tau\ge2$.
Write $d_T$ for the leaf metric: distinct leaves with last common
ancestor of depth $j$ have distance $r_j$.  At time $t$ there are $k$
labeled hidden members at leaves,
and a leaf $x_t$ known to $\mathcal F_t$ is required.  Assume that,
conditional on $\mathcal F_t$, at least one hidden member is at $x_t$
with probability one.  Let $N_v(t)$ be the number in the subtree of $v$,
and put $m_v(t)=\E[N_v(t)\mid\mathcal F_t]$.  The expected sum of
the labeled members' successive traveled distances is denoted by $P$;
it includes their movement from time zero.  All conditional distributions
below are available from the observed history.

\begin{theorem}[Dominated tree allocation]\label{thm:allocation}
There are absolute constants $C_0,C_1$ and $\tau_0$ with the following
property.  For $\tau\ge\tau_0$, the observed histories determine amounts
$q_v(t)$ satisfying
\begin{equation}\label{eq:allocation-constraints}
 q_{\mathrm{root}}=k,\qquad
 0\le q_v\le C_0m_v,\qquad
 \sum_{i\prec v}q_i\le q_v,\qquad q_{x_t}\ge1.
\end{equation}
Moreover, with $\ell=1+\log(k+1)$,
\begin{equation}\label{eq:allocation-bound}
 V:=\E\sum_{t\ge1}\sum_{v\ne\mathrm{root}}
       r_{|v|-1}|q_v(t)-q_v(t-1)|
 \le C_1\ell P+D_{\mathcal T}.
\end{equation}
Here sums stop at the chosen finite horizon.  The finite constant
$D_{\mathcal T}$ depends only on the fixed tree, radii, and $k$;
it can be chosen independently of the input law, horizon, and number of
histories.  No initial allocation is prescribed.
\end{theorem}

The proof occupies Sections~\ref{sec:allocation}--\ref{sec:allocation-finish}.

At a vertex $v$, the local task is to distribute its amount $q_v$ among
the children without violating their serving lower bounds.  We construct
child allowances $d_i$, then remove at least $(\sum_i d_i-q_v)_+$ from
them; additional removals leave mass parked at $v$.  The removal choices
must remain stable under posterior fluctuations.  Section~\ref{sec:proportions}
constructs these choices, and Section~\ref{sec:allocation-finish} proves
that the resulting allocations meet their lower bounds and move slowly.

This section supplies the rank lower bounds and the allowances.  It uses
coarse estimates to control changes of the removal rule's parameters,
and finer estimates for the final allowances.  Only the coarse variation
will be multiplied by logarithmic potential slopes; keeping these two
budgets separate preserves the single factor of $\ell$ in the theorem.

Rank coordinates for subtree occupancies also occur in the assignment
polytope of \cite[Section~3.2]{BCLLM18}. Here they are probabilities of
the hidden comparator's occupancies, conditional on the observed history.
For $1\le a\le k$, define the rank probabilities and their filtered values by
\[
 p_{va}(t)=\Pr(N_v(t)<a\mid\mathcal F_t),\qquad
 p^-_{va}(t)=\Pr(N_v(t-1)<a\mid\mathcal F_t).
\]
Thus $\E[p^-_{va}(t)\mid\mathcal F_{t-1}]=p_{va}(t-1)$.
The drift of this coordinate is
$\delta_t p_{va}=p_{va}(t)-p^-_{va}(t)$, rather than its entire change
between successive observations.  For the star at an internal vertex $v$,
put
\begin{equation}\label{eq:star-drift}
 E_v=\E\sum_{t\ge1}\left(
        \sum_{i\prec v}\sum_{a=1}^k|\delta_t p_{ia}|
             +\sum_{a=1}^k|\delta_t p_{va}|\right).
\end{equation}
These local budgets satisfy
\begin{equation}\label{eq:rank-drift-tree}
             \sum_{v\text{ internal}}r_{|v|}E_v\le6P.
\end{equation}
Indeed, for two integer counts $n,n'\in\{0,\ldots,k\}$,
\[
 \sum_{a=1}^k|\one_{\{n<a\}}-\one_{\{n'<a\}}|=|n-n'|.
\]
Use the coupling of successive hidden configurations and conditional
expectation to bound each rank drift by the corresponding indicator
change.  A nonroot vertex of depth $d$ receives weight at most
$r_{d-1}+r_d$ when all stars are summed.  If one labeled member moves
between leaves with last common ancestor of depth $j$, its indicators
change on the two paths strictly below that ancestor.  Its total charge
is at most
\[
 2\sum_{d=j+1}^J(r_{d-1}+r_d)
 \le\frac{2(1+\tau^{-1})}{1-\tau^{-1}}r_j\le6r_j.
\]
The root count is constant.  Summing over the members proves
\eqref{eq:rank-drift-tree}.

\subsection{Explicit rank functions and compatible lower bounds}

\begin{lemma}[Rank functions]\label{lem:rank-functions}
There are absolute constants $p_*,s_*,\epsilon_*>0$, with $p_*<s_*<1$,
and functions $l^\beta:[0,1]\to[0,1]$ for
$3\le\beta\le3+\epsilon_*$ such that:
\begin{enumerate}
\item $l^\beta$ is convex and nonincreasing, equals $1-\beta p$ on
      $[0,p_*]$, is positive on $[0,s_*)$, and is zero on $[s_*,1]$;
\item $l^\beta$ is affine and nonincreasing in $\beta$, and the square
      roots $\sqrt{l^\beta}$ have a common absolute Lipschitz constant;
\item if $g(p)=-\partial_\beta l^\beta(p)$, then
\begin{equation}\label{eq:rank-derivative}
 \frac1{96}l^\beta(p)\le g(p)\le l^\beta(p)
       \qquad(p\in[p_*/4,1]).
\end{equation}
\end{enumerate}
One may take $p_*=1/24$, $s_*=5/8$, and $\epsilon_*=1/100$.
\end{lemma}

\begin{proof}
Set $t_*=s_*-p_*=7/12$ and
\[
 A_*=\frac{18}{7},\qquad c_*'=\frac6{49},\qquad
 d_*'=\frac{1152}{343},\qquad b_*'=c_*'+d_*'t_*=\frac{102}{49}.
\]
For $0\le z\le1/100$ define
\begin{equation}\label{eq:explicit-rank}
 l^{3+z}(p)=
 \begin{cases}
  1-(3+z)p,&0\le p\le p_*,\\
  (s_*-p)^2\bigl[A_*-z(c_*'+d_*'(p-p_*))\bigr],
                                     &p_*\le p\le s_*,\\
  0,&s_*\le p\le1.
 \end{cases}
\end{equation}
The identities
\[
 A_*t_*^2=1-3p_*,\quad 2A_*t_*=3,\quad
 c_*'t_*^2=p_*,\quad -2c_*'t_*+d_*'t_*^2=1
\]
show that both values and first derivatives match at $p_*$.
They also show that $g$ is $p$ on the first interval and
$(s_*-p)^2(c_*'+d_*'(p-p_*))$ on the second.

Writing $x=s_*-p$, the middle expression for $l^{3+z}$ is
\[
            (A_*-zb_*')x^2+zd_*'x^3.
\]
Its two coefficients are nonnegative, with $A_*-zb_*'>2$.
Consequently it is positive before the cutoff, has nonpositive
$p$-derivative, and has positive second derivative.  Its value and
first derivative are zero at the cutoff.  Together with the matching
at $p_*$, this proves global convexity and monotonicity.
On the first interval the function is bounded away from zero, so its
square root has bounded derivative.  On the middle interval its square
root is
$x\sqrt{A_*-zb_*'+zd_*'x}$, which also has bounded derivative,
uniformly in $z$.  The matching and cutoff give the global Lipschitz bound.

For $p\in[p_*/4,p_*]$, the quotient
$g(p)/l^{3+z}(p)=p/(1-(3+z)p)$ lies between $1/96$ and $1$.
On the middle interval the quotient is
\[
 \frac{b_*'-d_*'x}{A_*-zb_*'+zd_*'x},\qquad 0<x\le t_*.
\]
It is at least $c_*'/A_*=1/21$ and at most
$b_*'/(A_*-(1/100)b_*')<1$.  Beyond the cutoff both functions vanish.
This proves \eqref{eq:rank-derivative}.
\end{proof}

Fix also the broader size function
\[
 l^{\mathrm{siz}}(p)=\frac{(3/4-p)_+^2}{(3/4)^2},\qquad
 L_v^\beta=\sum_{a=1}^k l^\beta(p_{va}),\qquad
 M_v=\sum_{a=1}^k l^{\mathrm{siz}}(p_{va}).
\]
Its square root is Lipschitz.  Both kinds of functions vanish at one
and have value one at zero.

The next comparison is a majorization argument in the sense of
Karamata \cite{Karamata32}. We include the convex-hinge proof to specify
the direction of the comparison and the role of the padded coordinates.

\begin{lemma}[Superadditivity]\label{lem:superadditivity}
For every internal vertex $v$ and every allowed $\beta$,
\begin{equation}\label{eq:rank-superadditivity}
 \sum_{i\prec v}L_i^\beta\le L_v^\beta,\qquad
 \sum_{i\prec v}M_i\le M_v\le m_v.
\end{equation}
No independence between child counts is required.
\end{lemma}

\begin{proof}
Suppose there are $d$ children.  Pool their $dk$ rank probabilities
and sort them increasingly, obtaining a vector $x$.
Let $y$ be the parent's $k$ increasingly ordered rank probabilities
followed by $(d-1)k$ ones.  Their total sums agree and equal $dk-m_v$.
If the first $n$ coordinates of $x$ contain $n_i$ coordinates from child
$i$, those coordinates can be chosen as the first $n_i$ of that child.
Since $N_v=\sum_iN_i$, conditional expectation gives
\[
 \sum_{b=1}^n x_b
 =\sum_i\E[(n_i-N_i)_+\mid\mathcal F_t]
 \ge\E[(n-N_v)_+\mid\mathcal F_t]
 =\sum_{b=1}^n y_b.
\]
The last identity also holds for $n>k$, precisely because of the padding.
For every real $b$, an increasing vector $z$ satisfies
\[
       \sum_a(b-z_a)_+
       =\max_{0\le n\le dk}\left(nb-\sum_{a=1}^n z_a\right).
\]
Thus the sum of every convex hinge $(b-p)_+$ on $x$ is no larger
than its sum on $y$.  Affine functions have equal sums.  Every convex
piecewise linear function is an affine function plus a nonnegative
linear combination of these hinges.  Approximating a continuous convex
function by its linear interpolants proves the same comparison for
that function.  Apply it to $l^\beta$ and $l^{\mathrm{siz}}$;
the padded ones contribute zero.  Finally convexity and the endpoint
values give $l^{\mathrm{siz}}(p)\le1-p$, and
$\sum_a(1-p_{va})=m_v$.
\end{proof}

From now on set $c_*=10^{-4}$; thus
$100c_*<l^{\mathrm{siz}}(s_*)=1/36$.  Choose an absolute $c_\beta>0$
so small that
$c_\beta(1+\log(1/c_*))\le\epsilon_*/2$, and define
\begin{equation}\label{eq:vertex-beta}
 \beta_v=3+\frac{c_\beta}{\ell}
                \log\frac{k}{\max(M_v,c_*)}.
\end{equation}
It belongs to $[3,3+\epsilon_*/2]$, because $M_v\le k$ and
$\log(k/c_*)/\ell\le1+\log(1/c_*)$.
The invariant propagated down the tree will be
\begin{equation}\label{eq:allocation-lower-bound}
                         q_v\ge L_v^{\beta_v}.
\end{equation}
Superadditivity gives compatible child lower bounds at the parent's
parameter $\beta_v$.  Since $M_i\le M_v$, the child's own parameter
satisfies $\beta_i\ge\beta_v$, and its required bound is no larger.
The resulting difference is the slack available for choosing slowly
varying child amounts.  We will estimate this difference in terms of
the relative child sizes.
At the root all ranks are zero, $M_{\mathrm{root}}=k$, and the lower
bound is $k$.  At a required leaf $p_{v1}=0$, so the lower bound is
at least one.

\subsection{A rank tracker with a bounded affine potential}

The input to a star may fluctuate substantially when information is
revealed.  The following tracker charges only drift, using signed
filtering changes before it decides whether to refresh.

\begin{lemma}[Rank tracking]\label{lem:rank-tracking}
Let $p_t,p_t^-\in[0,1]$ obey
$\E[p_t^-\mid\mathcal F_{t-1}]=p_{t-1}$, and let
$\mathcal E=\E\sum_t|p_t-p_t^-|$.
If $f\ge0$ has an absolutely bounded Lipschitz constant for $\sqrt f$,
then for $0<\xi<1/2$ there is a held estimate $\widehat f_t$ with
\begin{equation}\label{eq:rank-relative-tracker}
 (1+\xi)^{-1}f(p_t)\le\widehat f_t\le(1+\xi)f(p_t),\qquad
 \E\sum_t|\widehat f_t-\widehat f_{t-1}|
       \le\frac{C}{\xi}(\mathcal E+1).
\end{equation}
In particular zeros are tracked exactly.  A held estimate
$\widehat p_t$ refreshed when $|p_t-\widehat p_{t-1}|>\delta$
has expected number of refreshes at most
$\delta^{-2}(\mathcal E+1)$.
\end{lemma}

\begin{proof}
Initialize a reference $p^\circ=p_0$ and hold $f(p^\circ)$ until
the displayed factor comparison fails, then replace $p^\circ$ by
$p_t$.  With the reference held, use
\[
 \Phi(p;p^\circ)=p(1-p)+(p-p^\circ)^2
              =(1-2p^\circ)p+(p^\circ)^2.
\]
This is affine in $p$, lies in $[0,1]$, and has slope of absolute
value at most one.  At every step, first evaluate its change from
$p_{t-1}$ to $p_t^-$ and then from $p_t^-$ to $p_t$, holding the old
reference.  The former change has conditional expectation zero, and
the latter is at most $|p_t-p_t^-|$.  A reference reset then decreases
$\Phi$ by exactly $(p_t-p^\circ)^2$.  Telescoping gives
\begin{equation}\label{eq:rank-reset-credit}
        \E\sum_{\text{resets}}(p_t-p^\circ)^2\le\mathcal E+1.
\end{equation}

For two nonnegative values $a,b$ failing factor $1+\xi$ comparison,
\[
 (\sqrt a-\sqrt b)^2
 =|a-b|\frac{|\sqrt a-\sqrt b|}{\sqrt a+\sqrt b}
 \ge \frac{\xi}{5}|a-b|.
\]
The inequality includes the case in which exactly one value is zero.
The Lipschitz bound for $\sqrt f$ therefore implies
$(p_t-p^\circ)^2\ge c\xi|f(p_t)-f(p^\circ)|$ at every reset,
proving \eqref{eq:rank-relative-tracker}.  The same potential proves
the absolute-threshold claim, since every reset then releases more
than $\delta^2$.  Filtering was evaluated at every step before the
reset test; no expectation was restricted to a selected reset event.
\end{proof}

Apply the absolute tracker with threshold $p_*/8$ to every child rank.
A rank is called \emph{core} when its held estimate is at most $p_*/2$.
After the update, every core rank has $p\le5p_*/8$, every noncore
rank has $p>3p_*/8$, and every rank with $p<p_*/4$ is core.
Let $K_t$ count the changed core flags in this star at step $t$, and set
\begin{equation}\label{eq:core-data}
 n_i=\#\{a:\text{$a$ is core at $i$}\},\quad
 B_i=\sum_{a\text{ core at }i}p_{ia},\quad
 S=\sum_iB_i,\quad
 f_i^\beta=\sum_{a\text{ noncore at }i}l^\beta(p_{ia}).
\end{equation}
We call \(f_i^\beta\) the \emph{noncore}, or \emph{flex}, contribution.
The core ranks lie in the affine part of every rank function, so
\begin{equation}\label{eq:core-decomposition}
 L_i^\beta=n_i-\beta B_i+f_i^\beta,\qquad
       n_i+\sup_\beta f_i^\beta\le C M_i.
\end{equation}
For the second bound, a core rank contributes at least
$l^{\mathrm{siz}}(5p_*/8)>0$ to $M_i$.  Any positive $l^\beta(p)$
has $p<s_*$, where $l^{\mathrm{siz}}(p)\ge1/36$ and $l^\beta(p)\le1$.
These termwise bounds prove the assertion with an absolute constant.
In particular,
\begin{equation}\label{eq:inactive-rank-vanishing}
 M_i\le10c_*\quad\Longrightarrow\quad
 n_i=0\ \text{ and }\ L_i^\beta=f_i^\beta=0\ \text{for all $\beta$}.
\end{equation}
Also, if the old flags are held while the ranks are filtered and drifted,
the resulting old-flag sum always satisfies $B_i\le n_i^{\mathrm{old}}$,
even if those old flags no longer describe small current probabilities.
This elementary bound will be used at wholesale changes.

\subsection{Coarse data, active epochs, and parameter gaps}

We now describe a schedule local to a star.  Its estimates are computed
from the final, post-filtering and post-drift, ranks on each step.
All elementary term trackers persist throughout the entire process,
including times when a child is inactive and times when the schedule
below is restarted.

Choose a sufficiently small absolute size tolerance $\delta_s>0$,
for example $10^{-3}$.  Track every term of $M_i$ by
Lemma~\ref{lem:rank-tracking} with this tolerance, and let
$\widehat M_i$ be their sum.  A child is either inactive, with $a_i=0$,
or active.  Turn an inactive child on when $\widehat M_i\ge5c_*$,
setting $a_i=\widehat M_i$.  Turn an active child off when
$\widehat M_i\le2c_*$.  Otherwise, during an active period hold $a_i$
until factor $1+\delta_s$ comparison with $\widehat M_i$ fails,
then replace $a_i$ by $\widehat M_i$.  Initialize with the same rules,
using either admissible state in the overlap.  These choices ensure
\begin{equation}\label{eq:active-size-accuracy}
 \begin{gathered}
 0.9M_i\le a_i\le1.1M_i\quad\text{at active children},\qquad
 M_i\ge c_*\quad\text{if active},\\
 M_i\ge10c_*\quad\Longrightarrow\quad\text{active}.
 \end{gathered}
\end{equation}
By \eqref{eq:inactive-rank-vanishing}, an inactive child has no core
ranks or positive rank lower bounds.

We also hold an accurate estimate of $X_v=\max(M_v,c_*)$.
For this purpose take a possibly smaller absolute tolerance $\delta_p$,
track the parent size terms, and clip their sum below at $c_*$,
obtaining $\widehat X_v$.  Hold $H_v$ equal to its last reference value,
refreshing upon failure of factor $1+\delta_p$ comparison.  Then
\[
 (1+\delta_p)^{-2}X_v\le H_v\le(1+\delta_p)^2X_v.
\]
Put
\begin{equation}\label{eq:held-beta-upper}
 X_v^{\mathrm{low}}=\max\{c_*,H_v/(1+\delta_p)^2\},\qquad
 \beta^+=3+\frac{c_\beta}{\ell}\log\frac{k}{X_v^{\mathrm{low}}}.
\end{equation}
Here $X_v^{\mathrm{low}}\le X_v\le k$, and
$X_v^{\mathrm{low}}\ge X_v/(1+\delta_p)^4$.  Thus, for any prescribed
small absolute $\epsilon>0$, choosing $\delta_p$ sufficiently small gives
\begin{equation}\label{eq:beta-upper-error}
       \beta_v\le\beta^+\le\beta_v+\epsilon/\ell.
\end{equation}
The parameter stays in the allowed interval and changes only when
$H_v$ is refreshed.

An \emph{epoch} starts with a \emph{wholesale change}, performed initially
and whenever $\beta^+$ changes or the cumulative $\ell_1$ variation of
the held size vector since the epoch started exceeds $A/100$.
Here $A=\sum_i a_i$ at the start is held for that epoch.  Empty to
nonempty or nonempty to empty changes always start a new epoch.
At a wholesale change, choose a dominant index $o$ if some
$a_o\ge0.85A$ and $A>0$.  It is unique.  Keep that mark until the
next wholesale change.  All other active indices, or all active indices
when there is no mark, are called regular.  With a dominant index the
regular indices are also called sides.  No elementary size or rank
tracker is restarted at an epoch boundary.

During a marked epoch, let $T_{\mathrm{side}}=\sum_{i\ne o}a_i$.
Set $U=T_{\mathrm{side}}$ at its start.  Subsequently hold $U$ until
factor $1+\delta_s$ comparison with $T_{\mathrm{side}}$ fails, then
refresh it to that side total.  This includes transitions to or from zero.
When sides exist set $u=U/A$; otherwise set $u=0$.  For active children
define
\begin{equation}\label{eq:gap-weights}
 h_i=1+\max\{0,\log(A/a_i)\}\quad(i\text{ regular}),\qquad
 h_o=u\quad(o\text{ dominant}).
\end{equation}
The empty active domain needs none of these parameters.

Finally prepare a coarse scale for a marked child.  Independently of
the fine estimates used below, track every child rank value
$l^3(p_{ia})$ with tolerance $\delta_s$, and write its estimate as
$\widehat l^{\mathrm{coarse}}_{ia}$.  With a dominant index set
\begin{equation}\label{eq:dominant-scale-definition}
             D=\sum_{a\text{ noncore at }o}
                       \widehat l^{\mathrm{coarse}}_{oa};
\end{equation}
set $D=0$ when no dominant index exists.  All these term trackers run
on every child and every rank even when their estimates are not currently
used in $D$.

Here is a precise inventory of the coarse charges.  Write
$a_i^-,a_i^+$ for the old and new held sizes on a step, and similarly
for other held quantities; these superscripts in this inventory refer
to step-end states, not to the filtered hidden configuration.
Let $\mathcal W_v$ be the wholesale steps after initialization,
$\mathcal H_v$ the steps at which $H_v$ is refreshed, and
$\mathcal U_v$ the non-wholesale steps at which $U$ is refreshed.
Define the nonnegative charge
\begin{equation}\label{eq:coarse-charge-definition}
\begin{aligned}
 C_{v,t}={}&K_t+
    \sum_{i:a_i^-\ne a_i^+}(a_i^-+a_i^+)
    +\one_{\{t\in\mathcal H_v\}}(H_v^-+H_v^+)\\
 &+\one_{\{t\in\mathcal W_v\}}\sum_i(a_i^-+a_i^+)
    +\one_{\{t\in\mathcal U_v\}}
                    (T_{\mathrm{side}}^-+T_{\mathrm{side}}^+)
    +|D^+-D^-|.
\end{aligned}
\end{equation}
The side totals in this expression use the unchanged dominant mark
on a non-wholesale step.  A wholesale change of $U$, including a
change of mark, is covered by the wholesale total in the preceding term.

\begin{lemma}[Coarse budgets and parameter gaps]\label{lem:parameter-gaps}
The data above satisfy, for an absolute constant $C$ and a finite
$D_v^{\mathrm{data}}$ depending only on the fixed star and $k$,
\begin{equation}\label{eq:coarse-data-budget}
             \E\sum_t C_{v,t}\le C E_v+D_v^{\mathrm{data}}.
\end{equation}
The constant is uniform in the law and horizon.  Every active child obeys
\begin{equation}\label{eq:parameter-gap}
        0\le h_i\le C\ell,\qquad
        \beta_i-\beta_v\ge\frac{c_g h_i}{\ell}
\end{equation}
for an absolute $c_g>0$.  If a dominant index and sides are present,
then $0<u<1/4$, and with $H=1+\log(1/u)$ every side satisfies
$h_i\ge cH$ for an absolute $c>0$.
In addition, $n_i+B_i+\sup_\beta f_i^\beta\le C a_i$ at active
children, and all these quantities vanish at inactive children.
\end{lemma}

\begin{proof}
We give the amortization details because paying a large size on every
small refresh would not suffice.  Lemma~\ref{lem:rank-tracking}, summed
over the fixed collection of size terms, bounds the expected total
variation of their estimates by $CE_v$ plus a fixed initialization
constant.  It bounds $\E\sum_tK_t$ in the same way.  Constants may depend
on the fixed absolute tolerances but never on $k$ or the degree.

For a scalar estimate $x$ held at a reference $b>0$, a factor
$1+\delta$ failure gives
\[
             b+x\le(1+2/\delta)|x-b|.
\]
The last difference is bounded by the variation of $x$ since its previous
reference setting; these intervals are disjoint.  This pays all active
to active changes of $a_i$ and all changes of $H_v$.
Turning a child on after an off period costs its new size, and is paid
by the variation from a previous value at most $2c_*$ to a value at least
$5c_*$.  The first such event may instead use an initialization charge.
For a turn-off, if the last held size exceeds $4c_*$, variation since
that setting pays it because the new estimate is at most $2c_*$.
Otherwise the charge is at most $4c_*$ and is paid by the traversal
from at least $5c_*$ at the start of this on period to at most $2c_*$
at its end.  Each on period and each off period is used only once in
these latter arguments.  They overlap the scalar refresh intervals
only a bounded number of times.  This proves the budget for all
changed-$a_i$ charges.  Clipping the parent sum below at $c_*$ does
not increase its variation, so the same argument applies to $H_v$.

When $\beta^+$ changes, the old and new active totals are at most
$C(X_v^-+X_v^+)\le C(H_v^-+H_v^+)$, and are therefore paid by that
parent refresh.  For a wholesale event from the vector-variation test,
let $R$ be the cumulative variation since the previous epoch started
and let $A$ be that epoch's initial total.  The old total is at most
$1.01A$, the new total is at most $1.01A+R$, and $R>A/100$.
Their sum is consequently at most $203R$.  The intervals of variation
assigned to successive epochs are disjoint.  If $A=0$, the new total
is at most $R$, and the same conclusion holds.  Discarding unused
variation credit when a parent refresh intervenes cannot increase its
total use.  This proves the wholesale budget.

On a fixed marked epoch, variation of $T_{\mathrm{side}}$ is at most
$\sum_i|\Delta a_i|$.  The scalar relative-refresh argument pays every
$U$ refresh by variation since its previous setting.  The actual old
side total is within factor $1+\delta_s$ of that reference, so the
old plus new side totals have the same bound.  A zero crossing is
included: the variation then pays the positive reference or positive
new total.  Setting $U$ at a wholesale change is paid by the wholesale
total.  Thus all its asserted charges are covered.

At a non-wholesale step the dominant index is unchanged.  The variation
of \eqref{eq:dominant-scale-definition} is at most the summed variation
of its coarse term estimates plus $CK_t$: a changed flag inserts or
removes one term, each at most one.  At a wholesale step, its old and
new values are at most $C(a_o^-+a_o^+)$, interpreted as zero when that
mark is absent; this follows from the relative estimates and
\eqref{eq:core-decomposition}.  The wholesale totals pay these changes.
This proves \eqref{eq:coarse-data-budget}.

For the gaps, at every state of an epoch with $A>0$ the current active
total is in $[0.99A,1.01A]$.  In an unmarked epoch every child has
$a_i\le0.86A$.  In a marked epoch the marked size is at least $0.84A$,
and the side total is at most $0.16A$; in particular the mark stays
active.  Therefore every regular child has $a_i\le0.86A$.
By \eqref{eq:active-size-accuracy} and superadditivity,
\begin{equation}\label{eq:size-log-comparison}
 \frac{M_v}{M_i}\ge
 1+\frac9{11}\frac{\sum_{j\ne i,\ a_j>0}a_j}{a_i}
 \ge1+\frac9{11}\left(\frac{0.99A}{a_i}-1\right).
\end{equation}
For $x=A/a_i\ge1/0.86$, the logarithm of the last expression is
at least $c(1+\log x)$ for an absolute $c>0$.  To see the uniformity,
the ratio to $1+\log x$ is continuous and positive on this interval
and tends to one as $x\to\infty$; on a sufficiently large tail it is
at least $1/2$, and on the remaining compact interval it has a
positive minimum.  This proves the regular comparison.
For the dominant child, use the first inequality in
\eqref{eq:size-log-comparison}, the comparisons
$T_{\mathrm{side}}\asymp U$, and $a_o\le1.01A$ to obtain
\[
                 \log(M_v/M_o)\ge cU/A=cu.
\]
Here $u\le(1+\delta_s)0.16<1/4$, so $\log(1+cu)\ge c'u$.
If no sides exist, $u=0$ and the claimed inequality is automatic.

At every active child $M_i\ge c_*$ and $M_v\ge M_i$, so the clipping
in \eqref{eq:vertex-beta} is inactive at both vertices.  Hence
\[
               \beta_i-\beta_v
                  =\frac{c_\beta}{\ell}\log(M_v/M_i),
\]
which proves the gap in \eqref{eq:parameter-gap}.
Active held sizes have an absolute positive lower bound and
$A\le1.1k$, proving $h_i\le C\ell$.  On a side,
$a_i\le T_{\mathrm{side}}\le(1+\delta_s)uA$, whence
\[
 h_i\ge1+\log(1/u)-\log(1+\delta_s)\ge H/2.
\]
The final bounds on core and flex data follow from
\eqref{eq:core-decomposition}, $B_i\le n_i$, and
\eqref{eq:active-size-accuracy}.  Inactivity was addressed in
\eqref{eq:inactive-rank-vanishing}.
\end{proof}

\subsection{Fine caps and their separate budget}

The proportions constructed in the next section will use only the
coarse data above and the raw core sums $B_i$.  The following finer
estimates enter only the final allocations.  Keeping their variation
out of the proportion potentials is necessary for the single factor
of $\ell$ in Theorem~\ref{thm:allocation}.

The parameter gap at a regular child leaves room to track its flex
contribution between $f_i^{\beta_i}$ and $f_i^{\beta_v}$.  A marked
child can have an arbitrarily smaller gap as its side mass tends to
zero.  We therefore use a small upper approximation at that child;
Section~\ref{sec:proportions} explains how this asymmetry changes the
removal rule.

\begin{lemma}[Caps]\label{lem:caps}
Given a sufficiently small absolute $\epsilon_b>0$, the remaining
absolute tolerances can be chosen so that there are caps
$d_i=n_i+b_i$, zero on inactive children, with
\begin{equation}\label{eq:cap-feasibility}
 \begin{aligned}
 f_i^{\beta_i}\le b_i\le f_i^{\beta_v}
                                      &&(i\text{ regular}),\\
 T:=f_o^{\beta_v}\le b_o\le(1+\epsilon_b/\ell)T,
 \qquad c_DT\le D\le C_DT
                                      &&(o\text{ dominant}).
 \end{aligned}
\end{equation}
Here $c_D,C_D>0$ are absolute, fixed independently of
$\epsilon_b$ and the fine tolerances; in particular $D=0$ if and only
if $T=0$.  At every child $d_i\le C M_i\le C m_i$.
Their variation away from wholesale changes satisfies
\begin{equation}\label{eq:fine-cap-budget}
 \E\sum_{t\notin\mathcal W_v}\sum_i|d_i(t)-d_i(t-1)|
          \le C\ell E_v+D_v^{\mathrm{cap}},
\end{equation}
where $D_v^{\mathrm{cap}}$ is uniform in the law and horizon.
At wholesale changes the old plus new caps are bounded by $C$ times
the old plus new active totals.
\end{lemma}

\begin{proof}
For every rank, at both endpoint parameters $3$ and $3+\epsilon_*$,
run separate relative trackers with tolerance $\rho/\ell$, where
$\rho>0$ is a small absolute constant to be selected.
Interpolate their estimates by the convex coefficients expressing
$\beta^+$ between these endpoints, and sum over the current noncore
ranks at child $i$.  Call the result $\widetilde f_i$.
The affinity in $\beta$ and the relative endpoint estimates imply
\[
 (1+\rho/\ell)^{-1} f_i^{\beta^+}
       \le\widetilde f_i\le(1+\rho/\ell)f_i^{\beta^+}.
\]
Every noncore rank has $p>3p_*/8$, so
\eqref{eq:rank-derivative} applies to it.  Together with
\eqref{eq:beta-upper-error}, this yields
\begin{equation}\label{eq:flex-estimate-error}
       |\widetilde f_i-f_i^{\beta_v}|
          \le\frac{e}{\ell}f_i^{\beta_v},\qquad
                      e\le C(\rho+\epsilon).
\end{equation}
For example,
$f_i^{\beta_v}-f_i^{\beta^+}
=(\beta^+-\beta_v)\sum_{\mathrm{noncore}}g(p_{ia})
\le(\epsilon/\ell)f_i^{\beta_v}$.
If the true sum is zero then both endpoint estimates, and therefore
the interpolated sum, are exactly zero, since all the functions have
the same cutoff.

For a regular child, $h_i\ge1$ and the parameter gap gives
\begin{equation}\label{eq:regular-flex-slack}
 f_i^{\beta_i}
 =f_i^{\beta_v}-(\beta_i-\beta_v)
                  \sum_{\mathrm{noncore}}g(p_{ia})
 \le(1-\gamma/\ell)f_i^{\beta_v},
             \qquad \gamma=c_g/96>0.
\end{equation}
Choose an absolute
$0<\lambda_b\le\min\{\gamma/4,\epsilon_b/4,1/4\}$,
and then choose $\rho,\epsilon$ so small that the $e$ in
\eqref{eq:flex-estimate-error} is at most $\lambda_b/4$.
Define
\[
 b_i=(1-\lambda_b/\ell)\widetilde f_i\quad(i\text{ regular}),
 \qquad
 b_o=(1+\lambda_b/\ell)\widetilde f_o\quad(o\text{ dominant}).
\]
For a regular child the upper estimate is at most $f_i^{\beta_v}$,
because $e\le\lambda_b$; the lower estimate is at least
$(1-(\lambda_b+e)/\ell)f_i^{\beta_v}$, which dominates
\eqref{eq:regular-flex-slack}.  For a dominant child,
$(1+\lambda_b/\ell)(1-e/\ell)\ge1$, whereas
$(1+\lambda_b/\ell)(1+e/\ell)\le1+\epsilon_b/\ell$.
This proves the cap comparisons.

On the noncore range \eqref{eq:rank-derivative} also gives
\[
 l^{\beta_v}\le l^3
       =l^{\beta_v}+(\beta_v-3)g
       \le(1+\epsilon_*)l^{\beta_v}.
\]
The coarse tracker defining $D$ therefore satisfies
\[
 (1+\delta_s)^{-1}T\le D
       \le(1+\delta_s)(1+\epsilon_*)T.
\]
These choices give the claimed $c_D,C_D$ and show that their values
precede all fine choices.  The bound $d_i\le C M_i$ follows from
\eqref{eq:core-decomposition} and the relative estimates.

It remains to account for variation.  Lemma~\ref{lem:rank-tracking}
with tolerance $\rho/\ell$ pays the summed variation of all endpoint
estimates by $C\ell E_v$ plus a fixed constant.  On a non-wholesale
step, $\beta^+$ and the dominant mark are unchanged, so their
interpolation coefficients and the factors $1\pm\lambda_b/\ell$
are fixed.  Changes of the noncore sets cost at most $CK_t$, as does
the variation of $n_i$.  An inactive child has zero true endpoint
values and hence zero endpoint estimates; setting its cap to zero
therefore adds no variation beyond the same term and flag changes.
This proves \eqref{eq:fine-cap-budget}.  At a wholesale step,
$d_i^-+d_i^+\le C(a_i^-+a_i^+)$, including zero inactive values,
and summing proves the remaining statement.
\end{proof}

All the additive terms in this section arise from bounded potentials
and bounded initialization charges for a fixed number of rank trackers.
There are only $k$ ranks per vertex of the fixed finite tree.
The term trackers are never reinitialized when an epoch changes or a
tree label reappears.  Thus their total additive allowance is uniform
over every law and horizon.  The coarse allowance
\eqref{eq:coarse-data-budget} and the fine allowance
\eqref{eq:fine-cap-budget} are separate: only the former will be
multiplied by the $O(\ell)$ slopes of the next section's potentials.

The remaining construction starts from the caps $d_i$ and distributes
the required reduction $\bigl(\sum_i d_i-q_v\bigr)_+$ between children.
The star's parameters, flags, and proportion trackers can all be
updated from its own rank inputs before $q_v$ is known.  The actual
allocations are then propagated from the root downwards.
